\chapter{Review of Related Literature}
This chapter explores related studies and projects based on ride sharing. Discovering the challenges and successes of past ridesharing services provides proponents with a better understanding of what key features to prioritize and how best to implement those features in the mobile application. A synthesis of the literature review is presented to provide a deeper insight into the overall process of the project.

A Proof of Concept of a Carpooling Ridesharing System for the USC community by Pepito et al.~\cite{pepito2022} concluded that the mobile ridesharing application is feasible for academic institutions. However, they recommended that their mobile application, “Kapiot”, be optimized and made more efficient. They listed their system challenges, such as a lack of a communication system, a filter mechanism for the available drivers, and a dynamic routing system for the drivers. In addition to technical issues, the researchers were unable to properly establish the guidelines for the CES points. Even if the point system passed their test cases, the points earned by the driver would not be justifiable by the scheme established by the CES Office. Based on their system, the proponents incorporated a communication system and a dynamic routing system in the project. A filter mechanism was not added as the feature did not apply to the project. Communication between the proponents and the CES Coordinator has been established to ensure that the points earned through the mobile application were justified based on their criteria for the voluntary act of ridesharing~\cite{pepito2022}.

To adopt artificial intelligence into the ride-matching algorithm, the proponents reviewed how the energy of sound, which the literature calls arousal, is described and measured for music, for speech, and across the two, and how the kind of sound in a vehicle can be told apart on a phone (Sections~\ref{sec:lit-audio} to~\ref{sec:lit-ondevice}). Based on these studies, the proponents incorporated AI that rates the arousal of the sound in the car as Lively, Calm, or Quiet, and used it as a condition when matching drivers and passengers during inbound and outbound trips. The AI describes the sound in the car and does not judge the mood or the emotional state of anyone \cite{zhang2026generalisation,won2017adaptive}.

Masoud and Jayakrishnan’s A Real-Time Algorithm to Solve the Peer-to-Peer Ride-Matching Problem in a Flexible Ridesharing System utilized proximity-based matching through ellipses. When a driver is in spatial proximity to the passenger’s pick-up point, the algorithm matches the two. Furthermore, their system would dynamically match drivers and passengers based on real-time location data and travel time matrices. This is to ensure detours are minimized when the driver's mobile application approaches the pick-up point. Aside from using AI in the ride matching, the proponents opted for a fixed distance proximity from the driver to a passenger’s pick-up point. The algorithm also checked for the number of pasengers with the driver before executing the ride-matching algorithm. This is to ensure that the driver is notified on time, allowing them to think and respond accordingly to any passenger’s request on the road \cite{masoud2017real}.

On the other hand, Dynamic Carpooling in Urban Areas: Design and Experimentation with a Multi-Objective Route Matching Algorithm by Mallus et al. uses a route-based matching algorithm for outbound trips. In their system, the driver was given a pre-planned route based on the start and end destination. Only passengers whose drop-off point aligned with this route were considered for matching. This would ensure minimal detours for the driver as he drops off the passengers. The proponents adopted this approach to the project as it alleviated the need for drivers to adjust their routes in order to accommodate the passenger’s drop-off point\cite{mallus2017dynamic}.

Regarding the use of non-monetary incentives, Use of Rideshare Services to Increase Participant Recruitment and Retention in Research: Participant Perspectives by Leavens et al. looked at student-participants and their willingness to engage in the original study and future studies. The provision of a rideshare service led to high recruitment and maintained retention from the participants through all the visits. Despite the presence of monetary compensation, the students preferred to accept the rideshare service since it was cheaper than commercial rates. About the project, it would be difficult to monetize the rideshare services because student-drivers would need to apply for Provisional Authority with the Land Transportation Franchising and Regulatory Board (LTFRB). To ensure that student-drivers were motivated to do rideshare services, the proponents collaborated with the USC-CES Office, who require undergraduate students to reach a certain number of CES Points as a requirement for their graduation. The ridesharing mobile application would provide a unique opportunity to earn CES points throughout a student’s travel \cite{leavens2019use}.

To understand the value of ridesharing in the Southeast Asian (SEA) countries, Governance of the Risks of Ridesharing in Southeast Asia: An In-Depth Analysis by Icasiano and Taeihagh highlights both the benefits and the risks of ridesharing as a sustainable transportation in SEA smart cities. The main benefits include a reduction of car ownership and emissions and improved transportation access. The risks include passenger safety, liability, and data privacy, and the authors conclude that these risks should be adequately governed so that people are not excluded from the benefits (see Section~\ref{sec:lit-risks}). The proponents therefore treated safety and privacy as design requirements of the application and of the testing, and not as secondary issues~\cite{icasiano2021governance}.

\section{Studies on Arousal in Music, Speech, and Sound Environments}
\label{sec:lit-audio}
The mood AI of SanCarGo describes the energy of the sound inside a vehicle. This section reviews why energy, which the literature calls arousal, was chosen as the axis of the mood, and what studies report about measuring it for music, for speech, and across the two.

Russell (1980)~\cite{russell1980circumplex} proposed the circumplex model of affect, in which affect sits on two axes: valence, which is how pleasant, and arousal, which is how energetic. Russell and Pratt (1980)~\cite{russell1980environments} showed that the same two axes describe the affective quality that people attribute to environments, and not only to the feelings of a person. Axelsson, Nilsson, and Berglund (2010)~\cite{axelsson2010principal} found that pleasantness and eventfulness are the two main dimensions of how people perceive sound environments. Mitchell, Aletta, and Kang (2022)~\cite{mitchell2022soundscape} described how ISO/TS 12913-2:2018 rates a sound environment on eight perceptual attributes (pleasant, vibrant, eventful, chaotic, annoying, monotonous, uneventful, and calm), which form a circumplex of pleasantness and eventfulness. In that standard, a calm soundscape is both uneventful and pleasant. For SanCarGo, these studies support rating a sound environment on the energy axis. The Calm label of SanCarGo describes low energy and ignores how pleasant the sound is, which is a deliberate simplification, because the AI is meant to describe the sound in the car and not to judge it.

Papadakis et al. (2023)~\cite{papadakis2023city} used the eight attributes of ISO/TS 12913-2:2018 to compare how 90 residents of a city, a town, and a village in Greece perceived the same sounds, and found that where people live changed how they rated attributes such as chaotic, vibrant, monotonous, and uneventful. For SanCarGo, this is a reminder that the rating of a sound environment depends in part on the listener. Yin, Shao, Hao, and Lu (2024)~\cite{yin2024soundscape} mapped 279 evaluations of urban green spaces in Chengdu onto pleasantness and arousal, and found that two of the three types of setting gave relatively weak arousal, and that in one type only two sound attributes, calmness and appropriateness, were related to arousal. Wu et al. (2025)~\cite{wu2025untranslatability} built a Chinese scale from the ratings of 264 participants on 132 sound excerpts and found that the translated global scale overestimated positive emotion and activation. For SanCarGo, these two are cautions: the arousal of a sound environment can be weak in ratings, and a scale that is translated can be biased. They are one more reason why the rating check is read as a sanity check only.

Aljanaki, Yang, and Soleymani (2017)~\cite{aljanaki2017deam} built the DEAM benchmark, a set of music excerpts that people rated for arousal and valence, as a basis for music emotion recognition. The agreement among the raters of arousal was moderate, with alpha values from 0.28 to 0.66. Zhang et al. (2018)~\cite{zhang2018pmemo} built the PMEmo dataset, which holds chorus clips of popular songs with arousal and valence annotations, and it is a second music set with its own raters. Sun et al. (2023)~\cite{sun2023segmented} studied music emotion recognition with segments of 1 to 5 seconds and classified arousal best in that range, with 81\% to 87\% accuracy in separating high from low arousal on PMEmo and on the 4Q data set. For SanCarGo, DEAM is the training source of the music head, PMEmo is kept for evaluation only, and the 1 to 5 second range is the reference for the length of the window (Section~\ref{sec:ai-window}). The moderate agreement of the raters also means that no model can be expected to match human ratings better than the raters match each other. Kang and Herremans (2025)~\cite{kang2025unified} trained one music emotion model on several datasets that use different kinds of labels, among them DEAM and PMEmo with arousal and valence, and reported that adding the datasets together improved results on single datasets. This preprint is cited for the use of human arousal ratings of music as the training signal, as SanCarGo does.

Zhang and Yao (2026)~\cite{zhang2026generalisation}, in a preprint, tested whether arousal models carry over between kinds of sound. A model transferred well within one kind of sound, with a Spearman correlation of about 0.59, but largely failed across kinds of sound, with a correlation of about 0.05. Ching and Widmer (2025)~\cite{ching2025gap}, in a preprint, trained on five music datasets with arousal and valence ratings (EmoMusic, DEAM, PMEmo, WTC, and WCMED) and showed that models generalise poorly to data of another distribution, and that training on several diverse datasets improved generalisation across datasets. This is one reason that the music head is tested on PMEmo, which it was not trained on. Coutinho, Deng, and Schuller (2014)~\cite{coutinho2014transfer} found that arousal carries across speech and music: the transfer from speech to music accounted for 70\% of the variance and the transfer from music to speech for 28\%, and arousal gained the most from learning shared across the two. Simonetta, Certo, and Ntalampiras (2024)~\cite{simonetta2024joint} also learned the emotions of music and of other sounds together. For SanCarGo, the first study is the reason for using a separate arousal head for each kind of sound and not one model for all sound, and the others support placing speech and music on the same arousal axis, so that a car playing a podcast receives a label in the same way as a car playing music.

Zhou, Zhang, and Duan (2023)~\cite{zhou2023arousal}, in a preprint, learned an arousal value for speech from the labels of emotion categories. Each category is given a fixed anchor on the arousal scale: angry 0.59, fearful 0.60, happy 0.51, disgust 0.35, sad $-0.27$, surprised 0.67, and neutral 0. Their model reached a concordance correlation coefficient (CCC) of about 0.63 to 0.67 on the IEMOCAP data set. Russell and Mehrabian (1977)~\cite{russell1977threefactor} had described emotion along three dimensions, pleasure, arousal, and dominance, which is the earlier basis for giving each emotion category a position on an arousal scale. Antoniou et al. (2023)~\cite{antoniou2023speech} reviewed the way speech emotion systems are evaluated and recommended that the test speakers must not appear in the training data, with leave-one-speaker-out validation as the practice to follow. For SanCarGo, the anchors of Zhou et al. are the target of the speech head, and the speech head is validated by leaving one actor out at a time (Section~\ref{sec:ai-eval}).

\section{Studies on Mood Categories and Listening Context}
\label{sec:lit-music}
Panda, Malheiro, and Paiva (2023)~\cite{panda2023audio} surveyed the audio features used for music emotion recognition. Music emotion is commonly described along an arousal dimension, which is the energy of the music, and a valence dimension, which is how pleasant it is. For SanCarGo, this supports describing the moods by the energy of the sound alone, so that Quiet, Calm, and Lively follow the arousal dimension.

Hu, Downie, Laurier, Bay, and Ehmann (2008)~\cite{hu2008mirex} reported the lessons of the 2007 Music Information Retrieval Evaluation eXchange (MIREX) audio mood classification task, which grouped moods into five clusters. In the human assessment that built the ground truth, the assessors did not reach agreement on 32.2\% (278 of 864) of the clips that were judged by two or three people, and agreement was lowest for the first cluster (57.6\%). Hu, Bay, and Downie (2007)~\cite{hu2007simplified} built a simplified ground-truth set for music mood classification, and clustering the mood tags gave three clusters as the best fit, judged by the silhouette value. G{\'o}mez Ca{\~n}{\'o}n, Herrera, G{\'o}mez, and Cano (2019)~\cite{gomez2019emotions} studied how the language of listeners and their comprehension of lyrics affect the emotions that they perceive in music. Lyrics comprehension raised agreement for some emotions and lowered it for others, and agreement on complex emotions was low, at about 0.2. The authors argue that low agreement in the annotations limits the performance of music emotion recognition models and their generalization to new data. For SanCarGo, these studies support using a small set of three moods, because people agree less as the categories become finer, and they show that the agreement between people sets a ceiling on how well any model can match human ratings.

North, Hargreaves, and Hargreaves (2004)~\cite{north2004uses} studied how people use music in everyday life and found that reactions to music depend on whether the listener chose it. For SanCarGo, this supports letting each passenger choose the moods that they accept, instead of assigning a mood to the ride. Karageorghis et al. (2021)~\cite{karageorghis2021participant} studied 27 young men in simulated urban driving and found that slow-tempo music chosen by the participant, and music chosen by the researchers, gave a better affective state for urban driving than traffic noise alone, which raised mental demand. Wen et al. (2019)~\cite{wen2019music} used a driving simulator with 20 participants and found that rock music gave a higher arousal level, a higher perceived mental workload, more variation of speed, and more lane crossings than light music or no music. Ghojazadeh et al. (2023)~\cite{ghojazadeh2023music} reviewed 19 studies and found that listening to music increased arousal and mental load, with effects on driving that were mixed and in some indicators adverse. Dibben and Williamson (2007)~\cite{dibben2007invehicle} surveyed 1,780 British drivers and found that about two thirds listened to recorded music or music radio while driving, and that music was reported as less distracting than conversation. For SanCarGo, these four studies are cited only for the link between the music in a car and the state and driving of the driver, and not for any threshold.

In ride-hailing, Constine~\cite{constine2019uber} reported in 2019 that Uber introduced a Quiet Driver Mode for Uber Black rides, in which the rider states before the request that a quiet ride with little conversation is preferred. This is an industry precedent for a rider stating a ride preference, but it concerns conversation, the driver is not required to follow it, and nothing detects the sound.

Roth (2008)~\cite{roth2008market} described what market design has learned about how markets work. One principle is thickness: a market works well when enough participants are present at the same time to find a match. For SanCarGo, this explains why each additional mood reduces the number of available drivers in each mood, and so the chance that a passenger finds a match (Table~\ref{tab:ai-chance}).

\section{Studies on Sound-Type Detection, Mixed Audio, and On-Device Classification}
\label{sec:lit-ondevice}
Won, Alsaadan, and Eun (2017)~\cite{won2017adaptive} classified the audio captured by a smartphone in a moving vehicle into music, speech, speech with music, and noise. The average accuracy across four road conditions was 93.3\% to 96.6\%, and the accuracy of the noise class was 98.4\% to 99.8\%. For SanCarGo, this shows that the kind of sound in a vehicle can be told apart on a phone, and its range is the published comparison for the sound-type step (Section~\ref{sec:ai-results}). Its data are not the data of SanCarGo, so it is a reference and not a target.

Phukan, Buduru, and Sharma (2023)~\cite{phukan2023embeddings} compared pretrained embeddings for speech emotion and found that speaker-recognition embeddings, namely x-vector and ECAPA, beat YAMNet for emotion categories. Niizumi et al. (2023)~\cite{niizumi2023byola} and Koh and Dubnov (2021)~\cite{koh2021embeddings} studied pretrained audio embeddings that stay frozen. Together they report that a linear head on a frozen embedding is the standard check of how good an embedding is, and that pooling the embeddings over time with the mean and the maximum together works well. For SanCarGo, these studies are the basis for a linear head on a frozen embedding, the x-vector for speech, and the mean and maximum pooling of the YAMNet embeddings for music. Cances, Guyot, and Pellegrini (2019)~\cite{cances2019postprocessing} evaluated post-processing for sound event detection. They tuned thresholds for each class on validation data, and they stopped unstable labels from flipping with a moving average and a two-threshold rule (hysteresis). For SanCarGo, this is the idea behind choosing the cutoffs of the sound-type step on tuning data and behind the confirm rule, which is a simple form of the two-threshold rule.

Batlle-Roca et al. (2023)~\cite{batlleroca2023background} measured when background music can be heard in broadcast television. Music was barely audible when it was 22.7 to 36.4 dB below the other sound, and inaudible below that. Venkatesh et al. (2021)~\cite{venkatesh2021synthesising} mixed speech and music to train a speech and music detector for radio broadcast. They measured loudness in loudness units (LU, ITU-R BS.1770-4) and report that there is no ideal loudness difference between speech and background music. The radio guidance that they cite recommends music at least 7 to 10 LU below the speech, and as much as 23 LU below when the music is very quiet. The range of 7 to 18 LU for their synthetic mixes was their own choice, made by trial and by listening, and it gave the best training set; it is not a measurement of real broadcasts. Liu and Elhilali (2025)~\cite{liu2025salience} found that listeners rank speech among the most important sounds in a complex scene, and Laugs (2020)~\cite{laugs2020mixed}, in a master's thesis, reported that models built for a single source fail on mixed-source audio. For SanCarGo, these studies show the levels at which a second source in a mix matters, and they give indirect support for checking speech before music when both are present. No study measures which source should be rated when both are clearly audible, so the order is a choice of the proponents (Section~\ref{sec:ai-window}).

Janko and Lu{\v{s}}trek (2021)~\cite{janko2021energy} chose the sleep lengths of a sensing system from a measured trade-off between accuracy and energy. Georgiev et al. (2014)~\cite{georgiev2014dspear} measured that keeping the processor awake to listen costs about 383~mW, and that co-processors are far cheaper but are not exposed by most phones. Lyu (2026)~\cite{lyu2026sections}, in a preprint, annotated the sections of 100 songs, and by the division of its figures by the proponents, a song section lasts about 25 seconds on average. For SanCarGo, the first study is the method for choosing an update interval from a measured trade-off, which the proponents applied with accuracy and lag because they did not measure energy (Section~\ref{sec:ai-results}). The second shows why continuous listening is costly, and the third is a reference for how long the sound of a car can stay the same.

Zhao et al.~\cite{zhao2022survey} surveyed deep learning on mobile devices, including its applications, optimizations, and challenges. They also cover training on the data that is held on the device. Lane and Georgiev (2015)~\cite{lane2015deep} asked whether deep learning can revolutionize mobile sensing, and built a deep neural network engine for the low-power processor of a phone, noting that its small program and memory space limit the model to carefully built simple networks. For SanCarGo, these studies support classifying the sound on the driver's phone with small models. Not storing or sending the recording, so that only the resulting mood label is saved, is SanCarGo's own design decision (Section~\ref{sec:ai-runtime}) and not a finding of these studies.

\section{Software Testing Standards}
\label{sec:standards}
The studies above show what to build, but not how to organize the evidence that it works. For that, the proponents refer to ISO/IEC/IEEE 29119, an internationally agreed family of software testing standards developed by ISO/IEC JTC~1/SC~7/WG~26. It replaces several older standards, including IEEE 829 on test documentation, IEEE 1008 on unit testing, and BS 7925-1 and -2~\cite{iso29119site}. The parts relevant to this project are:

\begin{itemize}
    \item \textbf{Part 1, Concepts and Definitions}: the common terminology and ideas of software testing, including how testing is organized into levels such as unit, system, and acceptance testing~\cite{iso29119p1}.
    \item \textbf{Part 2, Test Processes}: a process model for organizing testing at the organizational, management, and test-execution levels.
    \item \textbf{Part 3, Test Documentation}: the documents that testing produces, such as test plans, test case specifications, and test reports.
    \item \textbf{Part 4, Test Techniques}: methods for designing test cases, for example from the specified requirements of a system or from its internal structure.
    \item \textbf{Part 5, Keyword-Driven Testing}: a keyword-based approach to automated test scripts.
\end{itemize}

Further parts address particular settings, including agile projects (Part 6) and AI-based systems (Part 11)~\cite{iso29119site}. Part 11 is relevant here because the mood classifier is an AI component whose output affects ride matching.

The proponents consulted the public summaries of these parts. The standards themselves are sold by ISO, except Parts 1 and 11, which ISO makes freely available~\cite{iso29119site}. SanCarGo is not assessed for conformance with the standard. It is used as a reference framework for naming test levels, recording test cases, and deriving expected results, as applied in the testing plan (Section~\ref{sec:phases}).

\section[Safety, Liability, and Privacy Risks of Peer-to-Peer Ridesharing]{Safety, Liability, and Privacy Risks of Peer-to-Peer Ridesharing}\label{sec:lit-risks}
Ridesharing puts people in a vehicle with others they may not know, so the literature treats its risks as matters to be governed and not dismissed. Icasiano and Taeihagh~\cite{icasiano2021governance} reviewed ridesharing in Southeast Asia and reported that safety was the top concern in an Indonesian consumer survey, that drivers harassing female customers had been reported, and that current rules place liability mainly on drivers, while the share of liability between drivers and platforms was still debated in the Philippine Congress. They also described privacy risks, such as a data breach affecting 380,000 accounts of a ride-hailing company in Singapore and a Philippine ruling that barred a ride-hailing company from monitoring riders with video and audio equipment. These findings apply to a student ridesharing application, where student-drivers carry liability and the application handles location and audio-related data.

Road safety is a further concern. In the Philippines, road traffic deaths rose by 39\%, from 7,938 in 2011 to 11,096 in 2021 according to the Philippine Statistics Authority, as reported at the launch of the WHO Global Status Report on Road Safety 2023 in the Philippines~\cite{whoPhRoad2023}. The report names motorcyclists, pedestrians, and bicyclists as vulnerable road users, which matters because many student-drivers ride motorcycles. Driver distraction adds to the risk. Using naturalistic driving data from 3,542 drivers, Guo et al.~\cite{guo2017distraction} found that visual-manual tasks, such as those performed on a cell phone, affected drivers of all ages, and that distraction from secondary tasks was a consistently higher threat for drivers younger than 30, the age group of university students. This is also the reason why the application is designed to limit unnecessary prompts to drivers (Section~\ref{sec:matching}), and why the testing rules forbid drivers from handling the phone while the vehicle is moving (Section~\ref{sec:uat-safety}).

Taken together, these studies show that the benefits of ridesharing do not remove its risks, and that the risks to drivers, passengers, and other road users must be identified, reduced, and monitored before any real ride is made. The proponents applied this in the safety protocol and risk assessment of Section~\ref{sec:safety}.

\section{Synthesis}
Taken together, these studies describe the energy of sound on one arousal axis for music, for speech, and for sound environments~\cite{russell1980circumplex,russell1980environments,axelsson2010principal,coutinho2014transfer,simonetta2024joint}. They show that arousal can be predicted from human ratings of music~\cite{aljanaki2017deam,zhang2018pmemo,sun2023segmented}, that it can be learned for speech from emotion categories with fixed anchors~\cite{zhou2023arousal}, and that a model carries over within one kind of sound but largely fails across kinds of sound~\cite{zhang2026generalisation}. The kind of sound in a vehicle can be told apart on a phone~\cite{won2017adaptive}, pretrained embeddings with a linear head are a standard and compact way to build the arousal models~\cite{phukan2023embeddings,niizumi2023byola,koh2021embeddings}, and thresholds and smoothing can be set on validation data~\cite{cances2019postprocessing}. They also share limitations: the agreement among human raters of arousal is only moderate~\cite{aljanaki2017deam,hu2008mirex,gomez2019emotions}, the audibility and loudness values for mixed audio come from television and radio and not from a car~\cite{batlleroca2023background,venkatesh2021synthesising}, and no study says which source should be rated when speech and music are both clearly audible. None of the studies matched drivers and passengers by the sound inside a vehicle, or tested such a classifier on phones in a ride. SanCarGo builds on this literature by rating the arousal of the sound in the car with a sound-type step and one head for each kind of sound, by stating the basis of each choice as a study, a range given by studies, or its own reasoned choice, and by reporting the limits of the AI in Section~\ref{sec:ai-limits}. For the testing of the whole application, the standards reviewed in Section~\ref{sec:standards} give the project a common vocabulary of test levels and test documentation. Studies of mobile deep learning support classifying sound on the phone with compact models~\cite{zhao2022survey,lane2015deep}, and studies of sensing energy support choosing the update interval from a measured trade-off~\cite{janko2021energy,georgiev2014dspear}.

The literature review provides a path for the proponents to develop and deploy a successful ridesharing mobile application. Understanding the past literature allows for flexible adaptation and creative innovation to build the same mobile application, yet designed to fit the modern technologies. By aligning system methods and processes with the project’s objectives, it was meant to ensure that the project attained its goal of improving student transportation and reducing environmental impact.

